\section{Scenario controls must preserve a fixed scientific workload}
\label{exec23:dss-scope}

The previous orchestration comparison motivates a closer question for a
nightly bank offer analysis: which requested effects survive a partly failed
scenario? Dataiku's detailed official manuals describe useful controls for
handling failures, testing a flow and initiating another run. Their current
technical text supports a more precise comparison than an index of scenario
features. The observations below concern the saved DSS 15 and DSS 14 manuals
acquired after the 7 October 2026 cutoff. They establish documented mechanisms,
with the historical qualification developed below; they do not establish an
executed adapter or an exhaustive account of the product.
\citep{exec23_dss_latest_flow,exec23_dss_latest_triggers,exec23_dss_latest_steps,
exec23_dss_dss14_flow}

\subsection{A tolerated operating error leaves a scientific question}

The step-control manual specifies sequential execution until the first failure
or completion of all steps. Some steps, including Build/train, can downgrade
an error to a warning and let later steps run. An example uses Ignore failure
when attempting to kill another scenario. Individual steps can be disabled
with Never, run with Always regardless of previous failures, or use conditions
on earlier failures, the current outcome or variables. The documented scenario
outcomes are SUCCESS, WARNING, FAILED and ABORTED. These distinctions are useful:
cleanup can still run after a failure, and a later diagnostic need not be lost.
They leave the application responsible for deciding whether its requested
causal effects have actually been returned.
\citep{exec23_dss_latest_flow,exec23_dss_dss14_flow}

Let $\mathcal Q$ be the declared set of cohort requests, $w_q$ their fixed
population weights with $\sum_{q\in\mathcal Q}w_q=1$, $\tau_q$ the cohort's
mean potential-outcome contrast between contact and no contact, and $c_q$
the incremental cost of contacting a member. The declared all-contact policy
has incremental net value, in dollars per member,
\begin{equation}
 V_{\mathcal Q}=\sum_{q\in\mathcal Q}w_q(\tau_q-c_q).
 \label{exec23:dss-target}
\end{equation}
For the following exact authored construction, cohort effects are stipulated
truth, not estimates whose identification or precision has been established.
There are two equal-weight cohorts, with $\tau_A=10$, $\tau_B=-30$ and
$c_A=c_B=1$. No contact has zero incremental value and cost. Thus
\begin{equation}
 V_{\mathcal Q}=\tfrac12(10-1)+\tfrac12(-30-1)=-11.
 \label{exec23:dss-full-net}
\end{equation}
If the effect for cohort B is missing after a failed or skipped step, averaging
the available cohort instead reports $V_{\rm available}=9$. The apparent
improvement of 20 dollars comes from changing the population. It reverses the
financial sign without any sampling noise. A tolerated error that leaves the
same effect missing creates the same arithmetic problem. These are constructed
missing-output illustrations; they are not recorded Dataiku run outcomes.
\citep{exec23_dss_denominator_reference}

The proposed scientific ledger keeps two requested effects and one returned
effect in every such illustration. It retains the failure, skip or tolerated
error as a distinct disposition, the declared weights and the missing result.
It can describe the available cohort's value for that cohort, but cannot
rename it the original population value. Assigning zero to a missing effect
would add a substantive assumption. Even renormalizing costs to the returned
cohort would change the all-contact policy being evaluated.

An explicit range can communicate a limited answer. If a separate substantive
argument justifies $\tau_B\in[-50,50]$, with cohort A's effect and the costs
still stipulated, the original target obeys
\begin{equation}
 V_{\mathcal Q}=4+\tfrac12\tau_B\in[-21,29].
 \label{exec23:dss-bounds}
\end{equation}
This identified set follows from assumed bounds, rather than statistical
confidence. It gives no guarantee of a positive net return. The route to a
sharper answer is to recover the missing effect, defend tighter bounds or
declare a different policy and population explicitly. Reporting only the
successful cohort cannot supply the lost information.
\citep{exec23_dss_denominator_reference}

\subsection{A trigger observes its rule, not a complete causal snapshot}

The trigger manual distinguishes initiation from execution. Several active
triggers are evaluated independently; any true condition can promote a scenario
for execution, which need not start immediately. For filesystem datasets,
change detection uses filenames, lengths and modification times, with an
optional marker file. A dataset-change trigger does not detect SQL data
changes; a SQL-query trigger compares the query's output. The suggested row
count or maximum-column query is a useful proxy, whose adequacy depends on the
scientific request.~\citep{exec23_dss_latest_triggers}

For example, two fixed records can retain both row count and maximum date
while cohort B's stored contrast changes from $-30$ to $30$. In the same
stipulated equal-weight population, net value then changes from $-11$ to $19$.
A count-and-maximum-date proxy cannot distinguish these two tables. This is a
logical counterexample for that chosen query, rather than a claim that SQL
triggers cannot use a stronger query or that the platform loses changes.
The proposed adapter should bind the relevant data snapshot and information
clock, with a change rule appropriate to those dependencies.
\citep{exec23_dss_denominator_reference}

By default, Suppress triggers while running is enabled: a second trigger can
go unevaluated during a long run. If suppression is disabled, an evaluation
finishing while the scenario is still running is ignored. For eligible tracked
events, the manual describes disabling suppression, enabling Delay runs and
disabling Squash delayed runs to retain a run per trigger. It also warns that
serial scenarios accumulate delayed runs if arrival throughput exceeds service
throughput. Squashing limits the delayed queue to at most one run and can
cancel an intermediate run. Time-based, manual and hidden Run another scenario
triggers are excluded from run delaying.
\citep{exec23_dss_latest_triggers}

These controls support different operational objectives. A latest-snapshot
analysis may deliberately combine changes, whereas a predeclared event study
may need every eligible event and its clock. The latter requires a durable
event/request record and reconciliation against the declared workload.
Trigger counts alone cannot certify that denominator, and every raw data
change need not be a scientific event. The application must define the
relationship before choosing suppression or squashing. This is a proposed
contract to qualify, rather than demonstrated event completeness.

\subsection{Existing tests help define a bounded integration experiment}

The scenario-step manual documents checks and data-quality rules: the step
fails when at least one selected check fails, and rule outcomes become
variables for later steps. It also documents project pytest selection and a
flow integration test that swaps in stable reference inputs, rebuilds outputs
and compares them with reference datasets. These mechanisms could support a
prototype comparison using public or generated data. A causal extension would
need reference targets, weights, information clocks, refusals and uncertainty
alongside ordinary output comparisons; the documentation does not show that
such an extension has been executed.
\citep{exec23_dss_latest_steps}

Run another scenario waits for completion and takes the child's outcome.
The manual says the initiating step's user is used regardless of the child
scenario's Run as setting. A proposed adapter must therefore preserve the
observed run identity together with the scientific request, then test the
relevant configuration. This narrow documented behavior does not establish
the whole authorization model. Similarly, a Clear step deletes dataset or
folder contents; persistent scientific evidence must have its own deliberate
retention arrangement.~\citep{exec23_dss_latest_steps}

The economical next experiment is a small fixed manifest containing complete,
failed, skipped and tolerated-without-effect requests. Check the returned
scientific ledger, weights and source identities before measuring retries or
runtime. Pair this with two data versions that share a selected trigger proxy
but differ in a scientific dependency. A successful native experiment would
support only its frozen release, configuration and manifest. Its merit should
be established with public/generated inputs before purchasing a large bank
study. No native execution, licensed entitlement or measured commercial cost
is asserted here.

\subsection{The cutoff requires a patch identity, not just a major version}
\label{exec23:dss-cutoff}

The current release index lists DSS 15.0.0 on 14 August 2026 and DSS 15.0.2 on
24 September 2026, both before the client cutoff. It lists DSS 15.1.0 on
8 October and DSS 15.2.0 on 15 October, both later. DSS 15 therefore cannot be
dismissed as entirely post-cutoff. Equally, its current major-version manual
cannot silently stand for the exact 15.0.2 technical bytes. The DSS 14 release
index lists 14.7.5 on 24 September. Its inspected step-control body matches
the DSS 15 body, which supports a current textual comparison but does not
authenticate the historical patch revision.
\citep{exec23_dss_latest_releases,exec23_dss_dss14_releases,
exec23_dss_latest_flow,exec23_dss_dss14_flow}

Release-index headings for Agent Skills and MCP server do not establish their
detailed behavior. Those claims need the relevant primary technical text and
edition. The detailed scenario review narrows one earlier source omission;
exact cutoff originals for the selected commercial server manuals and complementary market coverage remain to be investigated. The following SAS, Domino and Posit client studies supply selected source contracts and preserve their distinct remaining server and deployment questions. Rights,
prices and native qualification remain separate questions. This comparison
supports reuse of documented orchestration mechanisms as candidates, while
requiring an explicit scientific workload before an integration is accepted.
